% ===========================================================================
\section{Alcance}
% ===========================================================================

Este informe documenta la Etapa 2 del compilador de \textsc{MiniJava}: el \textbf{Analizador
Sintáctico Descendente Recursivo Predictivo}. Su objeto central es la \textbf{gramática}: se
parte de la que publica la cátedra, se la diagnostica, se la corrige a LL(1) paso a
paso, se la extiende con los nueve logros perseguidos y se verifica el resultado.

El criterio de exposición es que \textbf{cada corrección quede localizada en un paso concreto}.
Por eso las transformaciones no se presentan como un salto de la gramática vieja a la nueva, sino
en cuatro pasos:

\begin{enumerate}
  \item la \textbf{regla original}, tal cual el documento de la cátedra;
  \item el \textbf{defecto} detectado, con los conjuntos de Primeros que se solapan;
  \item la \textbf{transformación} aplicada, citando el esquema de la teoría que la justifica;
  \item el \textbf{resultado} y la verificación de que los Primeros quedan disjuntos.
\end{enumerate}

\subsection{Qué está y qué no está en este informe}

Está la gramática y su tratamiento. No está la implementación en Java del analizador, que se
describe en la documentación técnica del repositorio; sí está, en cambio, todo lo que la
implementación necesita decidir antes de escribirse: qué producción elegir ante cada token, dónde
la gramática es ambigua y con qué criterio se resuelve, y qué construcciones quedan deliberadamente
fuera del alcance del análisis sintáctico.

\subsection{Qué chequea y qué no chequea un analizador sintáctico}
\label{sec:alcance-chequeos}

Conviene fijarlo desde el principio, porque delimita todo lo que sigue.

El analizador sintáctico verifica que la secuencia de tokens esté en un orden que la gramática
pueda generar: la estructura de clases, interfaces, miembros, bloques y sentencias; que los
delimitadores abran y cierren; que una expresión esté bien formada; que un método tenga cuerpo y
un método de interfaz no.

No verifica nada que dependa del \emph{significado} de los nombres. Todo lo siguiente atraviesa el
análisis sintáctico sin objeción y corresponde al analizador semántico:

\begin{center}
\begin{tabular}{@{}p{0.46\textwidth}p{0.44\textwidth}@{}}
\toprule
\textbf{Pasa el análisis sintáctico} & \textbf{Por qué es semántico} \\
\midrule
\texttt{int x = "hola";}                    & compatibilidad de tipos \\
\texttt{noDeclarada = 5;}                   & declaración y alcance de nombres \\
\texttt{5 = x;}                             & que el lado izquierdo sea asignable \\
\texttt{class A extends UnaInterfaz \{\}}   & que el padre sea realmente una clase \\
\texttt{m(1,2,3)} con \texttt{m} de un argumento & aridad contra la declaración \\
\texttt{int x; int x;}                      & declaraciones duplicadas \\
\texttt{return;} en un método \texttt{int}  & compatibilidad del valor de retorno \\
una lambda que usa una variable de afuera   & captura de variables \\
\bottomrule
\end{tabular}
\end{center}

El último caso merece atención porque afecta a un logro. El enunciado de \emph{MiniLambdas E2}
prohíbe la captura de variables locales y de parámetros, pero eso \textbf{no es decidible
sintácticamente}: el analizador no sabe qué nombres están en alcance en cada punto. Lo único de
ese logro que la gramática puede imponer es la \emph{forma} --- parámetros sin tipo, cuerpo de una
única expresión, sin llaves --- y así se hace. La restricción sobre la captura queda anotada como
responsabilidad del analizador semántico.
